\chapter{Software Implementation}
\label{ch:implementation}

\section{Repository Structure}

The project follows a modular directory structure with separate packages for simulation, backend, ML, optimization, dashboard, and tests. All Python code is PEP~8 compliant, type-hinted, and documented. Configuration is externalized to \texttt{.env} and \texttt{settings.yaml}.

\section{Backend}

The FastAPI backend provides:
\begin{itemize}[nosep]
    \item MQTT subscriber service for real-time data ingestion.
    \item REST API endpoints for querying energy flows, predictions, alerts, and performance metrics.
    \item Pydantic schema validation for all inputs and outputs.
    \item SQLAlchemy ORM with 10 database tables (devices, measurements, PV, battery, environment, energy flows, predictions, optimization results, alerts, daily performance).
\end{itemize}

The database defaults to SQLite for zero-configuration development and can be switched to PostgreSQL for production by changing a single environment variable.

\section{Dashboard}

The Streamlit dashboard comprises 8 pages:
\begin{enumerate}[nosep]
    \item \textbf{Overview:} Real-time gauges and today's power profile.
    \item \textbf{Energy Flow:} Daily energy breakdown and self-consumption trends.
    \item \textbf{Historical:} Date-range selectable time-series plots.
    \item \textbf{Consumption Forecast:} Actual vs.\ predicted with MAE, RMSE, $R^2$.
    \item \textbf{PV Forecast:} Same structure for PV production.
    \item \textbf{Optimization:} Day-selectable MILP vs.\ baseline comparison.
    \item \textbf{Alerts:} Filterable anomaly table and timeline.
    \item \textbf{Performance:} Multi-day baseline vs.\ optimized comparison with annual projections.
\end{enumerate}

\section{Testing}

The test suite contains 47 unit and integration tests organized into 5 test files covering energy balance, battery SOC model, data preprocessing, optimization constraints, and simulation output. All tests pass on both the development environment (this container) and the target Windows machine (Python~3.13).
